\section{Results}
\label{sec:results}

\subsection{Manipulation Check: Speed Reduction Confirmed}
\label{sec:res_manipulation}

Before analyzing genuinely behavioral outcomes, we confirm the governor's mechanical effect. GEE results for mean speed (manipulation check) show STD reduces speed by 1.94~km/h relative to TUB ($p < 0.001$) and ECO by 4.22~km/h ($p < 0.001$), controlling for age, time of day, road class, and trip distance. Within-subject paired comparisons (76,726 TUB-STD user pairs, 10,710 TUB-ECO pairs) confirm: TUB-to-ECO speed reduction of 3.67~km/h ($d = 1.63$, large). This is mechanical and expected---we report it to validate the treatment but do not claim it as a finding.

\subsection{Phase~I: Cross-Sectional Prediction (February--November 2023)}
\label{sec:res_prediction}

The concurrent availability of governed and ungoverned modes on identical hardware provides a natural behavioral laboratory. We analyze which outcomes are strongly mode-dependent---and therefore likely to respond to governance removal---and which are mode-invariant.

\subsubsection{Safety margins}

The ungoverned mode produces substantially more aggressive riding. GEE-estimated mode effects (Table~\ref{tab:gee_results}; Figure~\ref{fig:forest}) show that relative to TUB, STD reduces harsh acceleration by 46\% (Poisson coefficient $\beta = -0.61$, $p < 0.001$) and ECO by 94\% ($\beta = -2.86$, $p < 0.001$). The harsh deceleration pattern is similar: STD $\beta = -0.63$, ECO $\beta = -3.27$ (both $p < 0.001$). Speed CV shows a statistically significant but small mode effect ($\beta_{\text{STD}} = -0.022$, $\beta_{\text{ECO}} = -0.039$), suggesting that speed variability relative to mean speed is largely a rider characteristic rather than a mode effect.

Within-subject paired comparisons (same users across modes) provide stronger causal evidence. For TUB-versus-ECO pairs ($n = 10{,}710$), harsh deceleration shows Cohen's $d = 1.36$ (large), harsh acceleration $d = 1.05$ (large), and cruise fraction $d = -1.11$ (large, favoring ECO). Crucially, travel time shows $d = 0.04$ (negligible), meaning governed riders spend the same amount of time traveling---they simply ride more smoothly.

\subsubsection{Demand margins}

TUB trips are approximately 19\% longer (928~m vs.~777~m for STD, paired $d = 0.34$, small), but travel time is virtually identical ($d = 0.007$, negligible). This suggests TUB riders cover more distance in the same time rather than taking fundamentally different trips. At the cross-sectional level, no evidence suggests that governance suppresses demand.

\subsubsection{Riding dynamics}

Cruise fraction shows a clear gradient: 0.259 (TUB), 0.320 (STD), 0.417 (ECO). Governed riders maintain more stable speeds, with TUB-ECO paired $d = -1.11$ (large). Zero-speed fraction is similar across modes (0.155--0.167), indicating that stop-and-go behavior is primarily determined by traffic conditions and route characteristics rather than governance.

\subsubsection{Feature importance}

LightGBM models with SHAP analysis (Figure~\ref{fig:shap}) reveal distinct feature importance hierarchies across outcome categories. For harsh events (binary classification, AUC = 0.734), the TUB mode indicator dominates (mean $|\text{SHAP}| = 0.49$), followed by ECO ($0.32$) and trip length ($0.25$). For speed CV (regression, $R^2 = 0.46$), log-distance ($0.12$) and trip duration ($0.11$) dominate, with TUB contributing only $0.03$---confirming that speed variability is primarily a trip-characteristic effect, not a mode effect. For cruise fraction ($R^2 = 0.25$), the pattern mirrors speed CV. These results validate the cross-sectional finding: governance affects \textit{safety} outcomes (acceleration aggressiveness) far more than \textit{dynamics} outcomes (riding patterns).

\begin{table}[htbp]
\centering
\caption{GEE mode effect estimates for primary outcomes (200,000-trip subsample, TUB = reference).}
\label{tab:gee_results}
\small
\begin{tabular}{lrrrr}
\toprule
\textbf{Outcome} & \multicolumn{2}{c}{\textbf{STD vs TUB}} & \multicolumn{2}{c}{\textbf{ECO vs TUB}} \\
\cmidrule(lr){2-3} \cmidrule(lr){4-5}
 & Coef. & $p$ & Coef. & $p$ \\
\midrule
\multicolumn{5}{l}{\textit{Safety (Poisson GEE)}} \\
Harsh accel count & $-0.610$ & $<0.001$ & $-2.863$ & $<0.001$ \\
Harsh decel count & $-0.627$ & $<0.001$ & $-3.269$ & $<0.001$ \\
\addlinespace
\multicolumn{5}{l}{\textit{Safety (Gaussian GEE)}} \\
Speed CV & $-0.022$ & $<0.001$ & $-0.039$ & $<0.001$ \\
\addlinespace
\multicolumn{5}{l}{\textit{Dynamics (Gaussian GEE)}} \\
Cruise fraction & $+0.061$ & $<0.001$ & $+0.160$ & $<0.001$ \\
Zero-speed frac. & $-0.008$ & $<0.001$ & $-0.007$ & $<0.001$ \\
\addlinespace
\multicolumn{5}{l}{\textit{Manipulation check (Gaussian GEE)}} \\
Mean speed (km/h) & $-1.942$ & $<0.001$ & $-4.221$ & $<0.001$ \\
\bottomrule
\multicolumn{5}{l}{\footnotesize All models control for age, time of day, day type, city, road class, and log(distance).} \\
\multicolumn{5}{l}{\footnotesize Exchangeable correlation, clustered by user, sandwich standard errors.} \\
\end{tabular}
\end{table}

\subsubsection{Predictions from cross-sectional evidence}

Table~\ref{tab:predictions} formalizes six testable predictions derived from the cross-sectional patterns above. For each outcome, we state the cross-sectional evidence (mode effect size and direction) and the predicted causal response if TUB were removed. The logic is straightforward: outcomes that show strong mode dependence should respond causally to mode removal; outcomes that are mode-invariant should not.

\begin{table}[htbp]
\centering
\caption{Six predictions derived from cross-sectional behavioral profiling (Phase~I).}
\label{tab:predictions}
\small
\begin{tabular}{clp{4.2cm}p{4.2cm}}
\toprule
\textbf{\#} & \textbf{Outcome} & \textbf{Cross-sectional evidence} & \textbf{Prediction if TUB banned} \\
\midrule
P1 & Harsh accel & Strong mode effect ($\beta = -0.61$, $d = 1.05$) & Significant reduction \\
P2 & Harsh decel & Strong mode effect ($\beta = -0.63$, $d = 1.36$) & Significant reduction \\
P3 & Speed CV & Weak mode effect ($\beta = -0.02$, SHAP $= 0.03$) & No significant change \\
P4 & Cruise fraction & Moderate mode effect ($\beta = +0.06$) & No definitive change \\
P5 & Trip count/users & No cross-sectional mode dependence & No demand cost \\
P6 & Compensation & Within-subject same travel time ($d = 0.007$) & No behavioral compensation \\
\bottomrule
\multicolumn{4}{l}{\footnotesize Cross-sectional coefficients from GEE (Table~\ref{tab:gee_results}); Cohen's $d$ from within-subject pairs.} \\
\end{tabular}
\end{table}

\subsection{Phase~II: Causal Validation (December 2023 Natural Experiment)}
\label{sec:res_validation}

The December 2023 system-wide TUB ban provides an exogenous policy shock to validate each prediction. We evaluate each using multi-outcome difference-in-differences, dose-response analysis, placebo testing, and within-user compensation testing.

\subsubsection{Which margins respond to governor imposition?}
\label{sec:res_causal}

The TWFE DiD (Figure~\ref{fig:did_coef}; Table~\ref{tab:did_results}) validates Predictions P1--P5: safety margins respond causally, while demand and dynamics margins do not.

Harsh acceleration ($\beta = -0.156$, SE $= 0.046$, $p = 0.0006$) and harsh deceleration ($\beta = -0.190$, SE $= 0.046$, $p = 4.4 \times 10^{-5}$) both show significant reductions that survive Bonferroni correction ($\alpha_{\text{Bonf}} = 0.00625$). A city with 50\% pre-ban TUB share experienced a reduction of approximately 0.08 harsh acceleration events per trip and 0.10 harsh deceleration events per trip after the ban.

In contrast, speed CV ($\beta = -0.004$, $p = 0.91$), cruise fraction ($\beta = +0.043$, $p = 0.17$), and zero-speed fraction ($\beta = +0.006$, $p = 0.69$) are all statistically indistinguishable from zero. Demand outcomes---trip count ($\beta = 34{,}753$, $p = 0.34$) and active users ($\beta = 7{,}169$, $p = 0.19$)---also show null effects. The TUB ban did not reduce ridership.

This pattern confirms Predictions P1--P5: outcomes with strong cross-sectional mode dependence (harsh events) respond causally, while outcomes with weak or no mode dependence (speed CV, cruise fraction, demand) do not.

\subsubsection{Dose-response}

City-level dose-response analysis confirms a monotonic relationship. For harsh deceleration, cities with larger TUB share reductions experienced proportionally larger reductions in harsh events ($r = 0.621$, $p = 1.5 \times 10^{-6}$, slope $= 0.202$, 95\% CI $[0.146, 0.340]$, $n = 50$ cities). Harsh acceleration shows a similar pattern ($r = 0.561$, $p = 2.2 \times 10^{-5}$). For the manipulation check (mean speed), the dose-response is $r = 0.567$ ($p = 1.7 \times 10^{-5}$).

\subsubsection{Placebo test}

All six primary outcomes pass the placebo test using September 2023 as a fictitious treatment date. The largest placebo coefficient is for mean speed ($\beta_{\text{placebo}} = -0.59$, $p = 0.25$), far from significance. For harsh events, placebo $p$-values are 0.14 (acceleration) and 0.33 (deceleration). This confirms that the December effects are treatment-driven, not artifacts of seasonal trends.

\begin{table}[htbp]
\centering
\caption{Multi-outcome TWFE DiD results (Dec 2023 TUB ban, Bonferroni-corrected).}
\label{tab:did_results}
\small
\begin{tabular}{lrrrrl}
\toprule
\textbf{Outcome} & $\hat{\beta}$ & SE & $p$ & 95\% CI & \textbf{Verdict} \\
\midrule
\multicolumn{6}{l}{\textit{Safety}} \\
Harsh accel & $-0.156$ & 0.046 & 0.0006 & [$-0.245, -0.066$] & Significant$^\dagger$ \\
Harsh decel & $-0.190$ & 0.046 & $<$0.001 & [$-0.281, -0.099$] & Significant$^\dagger$ \\
Speed CV & $-0.004$ & 0.033 & 0.906 & [$-0.068, +0.060$] & Null \\
\addlinespace
\multicolumn{6}{l}{\textit{Demand}} \\
Trip count & $+34{,}753$ & 36,706 & 0.344 & [$-37{,}189, +106{,}696$] & Null \\
Active users & $+7{,}169$ & 5,419 & 0.186 & [$-3{,}452, +17{,}791$] & Null \\
\addlinespace
\multicolumn{6}{l}{\textit{Dynamics}} \\
Cruise fraction & $+0.043$ & 0.031 & 0.171 & [$-0.019, +0.105$] & Null \\
Zero-speed frac. & $+0.006$ & 0.015 & 0.687 & [$-0.023, +0.035$] & Null \\
\addlinespace
\multicolumn{6}{l}{\textit{Manipulation check}} \\
Mean speed & $-1.436$ & 0.571 & 0.012 & [$-2.555, -0.317$] & Sig.\ (not Bonf.) \\
\bottomrule
\multicolumn{6}{l}{\footnotesize $^\dagger$Significant after Bonferroni correction ($\alpha = 0.05/8 = 0.00625$).} \\
\multicolumn{6}{l}{\footnotesize City and month fixed effects; cluster-robust SEs. $N = 549$ city-month observations.} \\
\end{tabular}
\end{table}

\subsubsection{Compensation test (Prediction P6)}
\label{sec:res_compensation}

\paragraph{Aggregate vs.\ within-user effects.}

The STD/ECO-only DiD (restricting the panel to governed trips only) reveals significant \textit{increases} in mean speed ($\beta = +1.73$, $p < 0.001$) and harsh events ($\beta_{\text{accel}} = +0.12$, $\beta_{\text{decel}} = +0.22$, both $p < 0.001$) after the TUB ban. At face value, this appears to indicate behavioral compensation: riders riding more aggressively on governed modes. However, speed CV, cruise fraction, and zero-speed fraction remain null ($p > 0.4$).

The mode-switcher within-user DiD resolves this apparent contradiction (Table~\ref{tab:compensation}). Comparing 24,357 switchers (used TUB pre-ban, STD/ECO post-ban) with 79,700 never-TUB users, all Cohen's $d$ effect sizes are \textit{negligible}:

\begin{itemize}
    \item Mean speed: $d = +0.133$ (DiD = $+0.41$~km/h)
    \item Harsh acceleration: $d = +0.040$
    \item Harsh deceleration: $d = +0.062$
    \item Speed CV: $d = -0.036$
    \item Cruise fraction: $d = -0.045$
    \item Zero-speed fraction: $d = -0.039$
\end{itemize}

The maximum $|d|$ is 0.133 (mean speed), well below the 0.2 threshold for a small effect. The aggregate STD/ECO increases are entirely attributable to \textit{composition}: former TUB users entering the STD/ECO pool are inherently more aggressive riders (pre-existing gap: switchers ride 0.60~km/h faster in STD/ECO even before the ban). There is no behavioral compensation on any unconstrained margin.

\paragraph{Placebo validation.}

The user-level placebo test (August--September $\rightarrow$ October 2023, no treatment) confirms the design. All six outcomes show negligible effects (maximum $|d| = 0.053$), with the largest effect in the \textit{opposite} direction to the treatment period. This rules out differential seasonal trends between switchers and never-TUB users.

\begin{table}[htbp]
\centering
\caption{Compensation test: mode-switcher within-user DiD effect sizes.}
\label{tab:compensation}
\small
\begin{tabular}{lrrrl}
\toprule
\textbf{Outcome} & \textbf{DiD estimate} & \textbf{Cohen's} $d$ & $p$ & \textbf{Interpretation} \\
\midrule
Mean speed (km/h) & $+0.405$ & $+0.133$ & $<0.001$ & Negligible \\
Harsh accel count & $+0.019$ & $+0.040$ & $<0.001$ & Negligible \\
Harsh decel count & $+0.034$ & $+0.062$ & $<0.001$ & Negligible \\
Speed CV & $-0.009$ & $-0.036$ & $<0.001$ & Negligible \\
Cruise fraction & $-0.009$ & $-0.045$ & $<0.001$ & Negligible \\
Zero-speed frac. & $-0.005$ & $-0.039$ & $<0.001$ & Negligible \\
\bottomrule
\multicolumn{5}{l}{\footnotesize 24,357 switchers vs.\ 79,700 never-TUB users. $|d| < 0.2$ = negligible.} \\
\multicolumn{5}{l}{\footnotesize Statistically significant $p$-values reflect large $N$, not substantive effects.} \\
\end{tabular}
\end{table}

\subsubsection{Prediction--validation summary}

Table~\ref{tab:validation} maps each cross-sectional prediction to its causal validation result. All six predictions are confirmed: outcomes that showed strong mode dependence in Phase~I respond causally in Phase~II, while mode-invariant outcomes remain null. The concordance between cross-sectional profiling and causal identification demonstrates that observational behavioral analysis has genuine predictive power for policy outcomes.

\begin{table}[htbp]
\centering
\caption{Prediction--validation summary: cross-sectional predictions vs.\ causal results.}
\label{tab:validation}
\small
\begin{tabular}{clp{2.8cm}p{3.5cm}c}
\toprule
\textbf{\#} & \textbf{Prediction} & \textbf{DiD result} & \textbf{Key statistic} & \textbf{Verdict} \\
\midrule
P1 & Harsh accel $\downarrow$ & $\beta = -0.156$, $p = 0.0006$ & Bonferroni significant & Confirmed \\
P2 & Harsh decel $\downarrow$ & $\beta = -0.190$, $p < 0.001$ & Bonferroni significant & Confirmed \\
P3 & Speed CV unchanged & $\beta = -0.004$, $p = 0.91$ & Null & Confirmed \\
P4 & Cruise frac.\ unchanged & $\beta = +0.043$, $p = 0.17$ & Null & Confirmed \\
P5 & No demand cost & $\beta_{\text{trips}}$: $p = 0.34$; $\beta_{\text{users}}$: $p = 0.19$ & Both null & Confirmed \\
P6 & No compensation & All $|d| < 0.14$ & Max $d = 0.133$ & Confirmed \\
\bottomrule
\multicolumn{5}{l}{\footnotesize Predictions from Phase~I (Table~\ref{tab:predictions}); validation from Phase~II DiD and compensation test.} \\
\end{tabular}
\end{table}

\subsection{Behavioral Escalation via TUB Adoption}
\label{sec:res_escalation}

\subsubsection{Experience-mediation analysis}

TUB adoption increases monotonically with experience: from 23.0\% of trips at trip~1 to 69.5\% at trips~101--200 (Figure~\ref{fig:experience}). This adoption trajectory mediates the experience-outcome gradient differently across categories:

Safety outcomes are heavily mediated by TUB adoption. Harsh acceleration shows 71.9\% mediation (the experience-related increase in harsh events is largely explained by increasing TUB usage), mean speed 65.6\%, and harsh deceleration 59.1\%. In contrast, dynamics outcomes are less mediated: speed CV 26.8\%, cruise fraction 19.9\%, zero-speed fraction 17.5\%. The implication: experienced riders' higher harsh event rates are primarily a mode selection story---they choose TUB more---not an inherent change in riding aggressiveness within a given mode.

\subsubsection{Multi-endpoint survival analysis}

Survival analysis with multiple event types (Figure~\ref{fig:survival}; Table~\ref{tab:survival}) confirms the TUB-mediated escalation pathway. For first harsh event (93.9\% event rate among 127,764 users), TUB users reach the event at median 2 trips versus 3 trips for non-TUB users (log-rank $\chi^2 = 9{,}121$, $p < 0.001$). The Cox model yields HR$_{\text{tub\_frac}} = 2.13$: doubling a user's TUB fraction doubles their hazard of first harsh event.

For first speeding (66.1\% event rate), the TUB effect is dramatic: median time-to-event is 7 trips for TUB users versus infinity (never) for non-TUB users (log-rank $\chi^2 = 60{,}963$, $p < 0.001$; HR$_{\text{ever\_tub}} = 5.73$, $C = 0.819$).

Notably, first high-CV trip shows a \textit{reversed} pattern: TUB users reach this event \textit{later} (median 10 trips) than non-TUB users (median 5 trips), with HR$_{\text{tub\_frac}} = 0.80$ (protective). This is mechanically consistent: TUB mode's higher absolute speeds produce a lower CV (relative variability), so TUB users are less likely to have high-CV trips. This finding reinforces that speed CV is primarily a trip characteristic, not a safety indicator that governance can improve.

\begin{table}[htbp]
\centering
\caption{Multi-endpoint survival analysis results ($n = 127{,}764$ users).}
\label{tab:survival}
\small
\begin{tabular}{lrrrr}
\toprule
\textbf{Event type} & \textbf{Event rate} & \textbf{Median (TUB / non-TUB)} & \textbf{Log-rank $\chi^2$} & \textbf{Cox HR (tub frac)} \\
\midrule
First harsh event & 93.9\% & 2 / 3 trips & 9,121 & 2.13 \\
First high-CV trip & 93.0\% & 10 / 5 trips & 5,402 & 0.80$^*$ \\
First speeding & 66.1\% & 7 / $\infty$ trips & 60,963 & 8.13 \\
\bottomrule
\multicolumn{5}{l}{\footnotesize $^*$Protective: TUB users take longer to reach high-CV events (mechanical, not behavioral).} \\
\multicolumn{5}{l}{\footnotesize $\infty$ = median not reached (>50\% of non-TUB users never speed). All log-rank $p < 0.001$.} \\
\end{tabular}
\end{table}
